\documentclass[12pt, a4paper]{article}
% --- Packages ---
\usepackage{geometry}
\geometry{a4paper, top=3cm, bottom=2.5cm, left=2.5cm, right=2.5cm}
\usepackage{amsmath, amssymb}
\usepackage{graphicx}
\usepackage{fancyhdr}
\usepackage{setspace}
\usepackage{xcolor}
\usepackage{titlesec}
\usepackage{float} 
\usepackage{caption} 
\usepackage{tcolorbox} 
\usepackage{tikz} 
\usetikzlibrary{calc, shapes, arrows, positioning, fit, shadows, arrows.meta, backgrounds, patterns, decorations.pathreplacing}
% --- Persian Support ---
\usepackage{xepersian}
\settextfont[Scale=1.1]{Yas} 
\setlatintextfont[Scale=1.0]{Times New Roman}
\setdigitfont{Yas}
% تنظیم فاصله خطوط
\setstretch{1.5}
% --- Header & Footer Setup ---
\pagestyle{fancy}
\fancyhf{}
\fancyhead[R]{\textbf{یادداشت آموزشی}}
\fancyhead[L]{\textbf{نظریه بازی‌ها}}
\renewcommand{\headrulewidth}{0.5pt}
\begin{document}
% --- Title Page ---
\begin{titlepage}
    \centering
    {\large \textbf{دانشکده مهندسی کامپیوتر} \par}
    \vspace{2.5cm}
    {\Huge \textbf{پروژه دوم} \par}
    \vspace{2.5cm}
    {\Large محمد آرین ایروانی \par}
    \vspace{0.5cm}
    
    \vspace{3cm}
    {\Large درس نظریه بازی‌ها \par}
    \vspace{0.5cm}
    {\Large استاد: دکتر حسین‌قربان \par}
    \vspace{0cm}
    \vfill
    {\Large تابستان ۱۴۰۵ \par}
\end{titlepage}
% --- Table of Contents ---
\newpage
\tableofcontents
\newpage
% --- Document Content ---

\section{رأی‌دهنده}

قاعده‌ای که در این بخش استفاده می‌کنم همون $F_k^{AV}$ است (همون انتخاب $k$ کاندیدای پرامتیازتر، با شکست تساوی بر اساس کوچک‌ترین اندیس، طبق چیزی که در متن پروژه اومده). برای PAV و CC که بعداً لازمم میشن، چون قاعده‌ی شکست تساوی مشخص نشده، خودم فرض کردم که بین کمیته‌های هم‌امتیاز، همونی انتخاب میشه که به‌صورت لغت‌نامه‌ای کوچیک‌تره -- فقط برای اینکه مثال‌ها یکتا و قابل چک‌کردن باشن، وگرنه روی نتیجه‌ی کیفی تاثیری نداره.

\subsection{رأی صادقانه و تعادل نش}

\subsubsection*{پرسش ۱}

جواب اینه که بله، همیشه تعادل نشه. راستش اولش فکر می‌کردم باید یه شرط خاصی روی پروفایل بذارم که این درست باشه (مثلاً فاصله امتیازها یا یه چیزی تو این مایه‌ها)، ولی وقتی روش فکر کردم دیدم اصلاً نیازی نیست -- صداقت یه استراتژی برتر ضعیفه، یعنی مهم نیست بقیه چیکار می‌کنن، رأی راست‌گفتن همیشه حداقل به‌خوبی هر رأی دیگه‌ست.

استدلالش اینجوریه: امتیاز هر کاندیدا $c$ رو می‌تونیم بنویسیم $s_c = t_c + \mathbb{1}[c\in B_i]$ که $t_c$ امتیازیه که از بقیه رأی‌دهنده‌ها میاد (که به رأی خود $i$ ربطی نداره). نکته‌ی مهم اینه که امتیاز هر کاندیدا فقط به بیت مربوط به همون کاندیدا تو رأی $i$ بستگی داره، نه به بقیه‌ی رأی $i$.

حالا فرض کن $i$ رأیش رو دقیقاً روی یه کاندیدا $c$ عوض کنه (بقیه رو دست نزنه). دو حالت داره:
\begin{itemize}
\item اگه $c\in A_i$ بود ولی $i$ تاییدش رو برداره: امتیاز $c$ یکی کم میشه. چون قاعده‌ی top-$k$ نسبت به امتیاز یکنواخته (بالا بردن امتیاز فقط می‌تونه کمک کنه، نه ضرر)، این کار فقط ممکنه $c$ رو از کمیته بندازه بیرون، نه برعکس. اگه بیفته بیرون، یکی دیگه ($d$) جاش میاد. مطلوبیت $i$ عوض میشه به اندازه‌ی $-1 + \mathbb{1}[d\in A_i] \le 0$ -- یعنی یا بدتر میشه یا همون میمونه.
\item اگه $c\notin A_i$ بود و $i$ الکی تاییدش کنه: امتیاز $c$ زیاد میشه، ممکنه وارد کمیته بشه و یکی دیگه رو بندازه بیرون. اینجا هم مطلوبیت $i$ برابر $-\mathbb{1}[d\in A_i]\le 0$ میشه.
\end{itemize}

خب پس هر تغییر تک‌کاندیدایی از رأی صادقانه، مطلوبیت رو یا کم می‌کنه یا ثابت نگه می‌داره. حالا هر رأی دلخواه $B_i$ رو می‌شه با یه سری از همین تغییرات تک‌کاندیدایی به $A_i$ رسوند، و چون اثر هر تغییر مستقل از بقیه‌ی رأیه (طبق چیزی که بالا گفتم)، این استدلال کل مسیر رو پوشش میده. پس حرکت از $B_i$ به سمت $A_i$ همیشه (ضعیفاً) بهتر میشه:
\[
u_i\big(F_k^{AV}(A_i,B_{-i})\big) \ge u_i\big(F_k^{AV}(B_i,B_{-i})\big) \qquad \text{برای هر } B_{-i} \text{ و هر } B_i.
\]
چون این برای هر $B_{-i}$ برقراره، یعنی $A_i$ استراتژی برتر ضعیفه. و وقتی همه استراتژی برتر دارن، خب معلومه که تعادل نش هم هست -- هیچکس با تغییر یک‌طرفه بهتر نمیشه.

\subsubsection*{پرسش ۲}

طبق چیزی که بالا نشون دادم، صداقت برای \emph{هر} $B_{-i}$ بهترین پاسخه، نه فقط وقتی بقیه هم صادقن. پس بله، همیشه بهترین پاسخه (شاید تنها بهترین پاسخ نباشه -- ممکنه چند رأی دیگه هم به همون اندازه خوب باشن -- ولی هیچ‌کدوم بهتر نیستن). به همین خاطر لازم نبود دنبال نمونه‌ی نقض بگردم.

\subsubsection*{پرسش ۳}

سه تا چیزی که تو صورت سوال اشاره شده (حساسیت به تغییر یک نفر، فاصله امتیازی، و شکست تساوی) دقیقاً همون چیزاییه که تو اثبات بالا استفاده کردم:

\begin{itemize}
\item امتیاز هر کاندیدا فقط به تاییدِ همون کاندیدا وابسته‌ست، نه به بقیه‌ی رأی. یعنی هیچ رابطه‌ی پنهانی بین کاندیداهای مختلف تو رأی یه نفر نیست که بشه ازش سوءاستفاده کرد.
\item اگه کاندیدای موردعلاقه‌ی $i$ خیلی بالا یا خیلی پایین‌تر از خط $k$-ام باشه، رأی $i$ اصلاً روی نتیجه اثر نداره (pivotal نیست)، پس دروغ گفتن هم فایده‌ای نداره. وقتی رقابت نزدیکه هم -- طبق همون اثبات -- بهترین کاری که $i$ می‌تونه بکنه دقیقاً همون کاریه که صداقت انجام میده.
\item شکست تساوی فقط به اندیس کاندیدا بستگی داره، نه به اینکه کی به کی رأی داده. اگه اینجوری نبود (مثلاً اگه تساوی بر اساس اینکه رأی‌دهنده‌های کوچیک‌تر به کی رأی دادن شکسته می‌شد)، شاید می‌شد با ترکیب رأی‌ها روی نتیجه‌ی تساوی اثر گذاشت.
\end{itemize}

خلاصه بگم، این سه‌تا با هم باعث میشن $F_k^{AV}$ تو دسته‌ی قاعده‌های «جداپذیر» قرار بگیره -- امتیاز هر کاندیدا مستقل حساب میشه -- و همین جداپذیری دقیقاً چیزیه که صداقت رو استراتژی برتر می‌کنه.

\subsubsection*{پرسش ۴}

اینجا رفتم سراغ PAV، و برخلاف $F_k^{AV}$، اینجا صداقت اصلاً تضمین نیست. اولش با دست چندتا نمونه امتحان کردم و پیدا نمی‌کردم، بعد یه اسکریپت نوشتم که همه‌ی حالت‌های ممکن رو برای نمونه‌های کوچیک چک کنه (کدش پیوسته)، و همون اولین اجراش یه نمونه‌ی خیلی کوچیک پیدا کرد: $n=2$ رأی‌دهنده، $m=3$ کاندیدا $\{a,b,c\}$، $k=2$.

$A_1=\{a,b\}$ و $A_2=\{a,c\}$. اگه هر دو صادقانه رأی بدن، امتیازای PAV اینجوری میشه (با $H_1=1$, $H_2=1.5$):
\[
\{a,b\}: 1.5+1=2.5, \qquad \{a,c\}: 1+1.5=2.5, \qquad \{b,c\}: 1+1=2.
\]
تساوی بین $\{a,b\}$ و $\{a,c\}$، که با شکست تساوی لغت‌نامه‌ای، $W=\{a,b\}$ برنده میشه. رأی‌دهنده‌ی ۲ که واقعاً $\{a,c\}$ رو دوست داره، از این کمیته فقط $a$ رو میگیره یعنی مطلوبیتش ۱.

حالا اگه رأی‌دهنده‌ی ۲ به‌جای گفتن رأی واقعیش، فقط بگه $\{c\}$ (یعنی تاییدش از $a$ رو مخفی کنه با اینکه واقعاً دوستش داره)، امتیازها میشه:
\[
\{a,b\}: 1.5+0=1.5,\qquad \{a,c\}: 1+1=2,\qquad \{b,c\}: 1+1=2.
\]
حالا $\{a,c\}$ (با تساوی با $\{b,c\}$، ولی لغت‌نامه‌ای $\{a,c\}$ میبره) برنده میشه، و رأی‌دهنده‌ی ۲ حالا هر دو کاندیدای موردعلاقه‌شو داره، مطلوبیتش شده ۲. یعنی با دروغ گفتن بهتر شد!

اولش این برام عجیب بود -- چطور مخفی کردن یه چیزی که واقعاً میخوایش کمکت می‌کنه؟ فکر کنم دلیلش اینه: تو PAV، وزن‌های $H_r$ کاهشین ($H_2-H_1=0.5$ کمتر از $H_1-H_0=1$)، یعنی گرفتن هر دو کاندیدای موردعلاقه با هم، ارزشش کمتر از دوبرابر گرفتن یکیشونه. وقتی رأی‌دهنده‌ی ۲ تاییدش از $a$ رو قایم می‌کنه، امتیاز $\{a,b\}$ کاملاً از دست میده (چون دیگه $a$ رو تایید نکرده)، ولی امتیاز $\{a,c\}$ فقط از $H_2$ به $H_1$ افت می‌کنه (چون $c$ هنوز هست). این ضرر نامتقارن باعث میشه رتبه‌بندی برعکس بشه.

جالبه که این دقیقاً همون چیزیه که PAV رو برای نمایندگی گروهی خوب می‌کنه (بخش بعد میبینیم)، ولی همین ویژگی درِ سوءاستفاده رو هم باز می‌ذاره.

\subsubsection*{پرسش ۵}

برای $F_k^{AV}$: چون اثباتِ پرسش ۱ اصلاً به $k$ ربطی نداشت، صداقت برای \textbf{هر} $k$ استراتژی برتره. یعنی تغییر $k$ هیچ فرقی نمی‌کنه.

برای PAV کنجکاو شدم ببینم چی میشه، پس کد رو روی $m=4$ و همه‌ی $k$ها اجرا کردم:

\begin{center}
\begin{tabular}{c|c}
$k$ & نمونه‌ی دروغ‌گفتن پیدا شد؟ \\
\hline
$1$ & نه \\
$2$ & آره \\
$3$ & آره \\
$4$ (=m) & نه
\end{tabular}
\end{center}

این الگو با یه فکر ساده توضیح داده میشه: وقتی $k=1$، چون هر رأی‌دهنده حداکثر یک کاندیدا تو $W$ داره، $H_{|A_i\cap W|}$ فقط بین $H_0=0$ و $H_1=1$ در نوسانه -- یعنی PAV دقیقاً میشه همون $F_k^{AV}$ (وزن‌های کاهشی اصلاً وارد بازی نمیشن)، پس طبق پرسش ۱ صادقه. وقتی $k=m$ هم فقط یه کمیته وجود داره ($W=C$) و هیچ رأیی چیزی رو عوض نمی‌کنه. برای مقادیر وسط، جایی که یه رأی‌دهنده ممکنه چند کاندیدای موردعلاقه‌ی هم‌زمان تو رقابت داشته باشه، همون اثر کاهشیِ $H_r$ فعال میشه و امکان دروغ‌گفتن پیدا میشه.

\subsection{رأی‌دهنده در شرایط اطلاعات ناقص}

\subsubsection*{پرسش ۱}

فضای نوع رأی‌دهنده‌ی $i$ همون $\Theta_i=2^C$ است -- نوعش دقیقاً همون مجموعه‌ای‌ه که واقعاً می‌پسنده، $A_i\in\Theta_i$. رأی‌دهنده نوع خودش رو دقیق میدونه ولی نوع بقیه رو نه، فقط یه باور $\mu_i(\cdot\mid A_i)\in\Delta(\Theta_{-i})$ داره. این باور می‌تونه از رأی‌دهنده به رأی‌دهنده فرق کنه، مثلاً اگه سلیقه‌ها با هم همبستگی داشته باشن. چیزی که نمی‌دونه دقیقاً همون رأی‌های واقعی بقیه‌ست، و در نتیجه امتیاز دقیق هر کاندیدا قبل از رأی خودش.

\subsubsection*{پرسش ۲}

مطلوبیت انتظاری برای یه رأی دلخواه $B_i$ (با فرض اینکه بقیه طبق $\sigma_{-i}$ عمل می‌کنن):
\[
U_i(B_i\mid A_i;\sigma_{-i}) = \mathbb{E}_{A_{-i}\sim\mu_i(\cdot\mid A_i)}\Big[u_i\big(F(B_i,\sigma_{-i}(A_{-i}))\big)\Big].
\]
فرقش با بازی اطلاعات کامل اینه که اونجا فقط یه پروفایل خاص $B_{-i}$ رو بهینه می‌کردیم، اینجا باید میانگین وزن‌دار روی همه‌ی پروفایل‌های ممکن رو بهینه کنیم. یه رأی ممکنه تو بعضی حالتا عالی باشه و تو بعضیا بد، مهم اینه که مجموعش (با وزن باور) از همه بیشتر باشه.

ولی چون تو بخش قبل نشون دادم $A_i$ برای \emph{هر} $A_{-i}$ (نه فقط یه حالت خاص) حداقل به‌خوبی هر رأی دیگه‌ست، این برتری زیر امید ریاضی هم حفظ میشه -- پس صداقت تو بازی بیزی هم، برای هر باور و هر استراتژی بقیه، بهترین پاسخه.

\subsubsection*{پرسش ۳}

فرض کنیم بقیه صادقانه رأی میدن. طبق چیزی که گفتم، برای هر $A_{-i}$ داریم
\[
\Delta(A_{-i}) := u_i\big(F_K^{AV}(A_i,A_{-i})\big) - u_i\big(F_K^{AV}(B_i,A_{-i})\big) \ge 0.
\]
این یه متغیر تصادفی نامنفیه. مطلوبیت انتظاری $B_i$ برابر مطلوبیت انتظاری $A_i$ میشه اگه و فقط اگه امیدش صفر باشه، و یه متغیر نامنفی امید صفر داره اگه و فقط اگه با احتمال یک صفر باشه. پس:
\[
U_i(B_i\mid A_i;\sigma_{-i}) = U_i(A_i\mid A_i;\sigma_{-i}) \iff \mu_i\big(\{A_{-i}: \Delta(A_{-i})=0\}\mid A_i\big)=1.
\]

از اینجا نتیجه میشه که صداقت تنها بهترین پاسخه اگه و فقط اگه برای هر $B_i\ne A_i$، این تساوی با احتمال یک برقرار \emph{نباشه}، یعنی
\[
\Pr\big(u_i(F_K^{AV}(A_i,A_{-i})) > u_i(F_K^{AV}(B_i,A_{-i}))\big) > 0.
\]

این رویداد دقیقاً همون حالتیه که $i$ تأثیرگذاره -- یعنی رأیش واقعاً چیزی رو عوض می‌کنه به نفعش. پس صداقت به‌عنوان تنها بهترین پاسخ، دقیقاً وقتی تضمین میشه که تو حمایتِ باور $i$ (یعنی جاهایی که $\mu_i$ بهشون احتمال مثبت میده) همچین سناریوهایی وجود داشته باشه. اگه رأی‌دهنده مطمئن باشه رأیش هیچوقت تعیین‌کننده نیست، صداقت هنوز بهترین پاسخ میمونه (چون هیچوقت ضرر نمیکنه) ولی دیگه تنها بهترین پاسخ نیست.


\section{برگزارکننده}

مطلوبیت رأی‌دهنده‌ها همون $u_i(W)=|A_i\cap W|$ می‌مونه، ولی حالا سوال اینه که خودِ کمیته چقدر خوبه از دید جمعی. برای PAV و CC همون قرارداد شکست تساوی لغت‌نامه‌ای بخش قبل رو نگه می‌دارم.

\subsection{نمایندگی گروهی}

\subsubsection*{پرسش ۱}

اول باید ببینیم این چهارتا (JR, PJR, EJR, Core) چه ربطی به هم دارن. جوابش اینه که:
\[
\text{Core-stable} \Rightarrow \text{EJR} \Rightarrow \text{PJR} \Rightarrow \text{JR}.
\]

PJR$\Rightarrow$JR راحته: JR دقیقاً حالت $\ell=1$ از PJR‌ه. وقتی $\ell=1$، شرط PJR میشه $|(\bigcup_{i\in S}A_i)\cap W|\ge 1$ که یعنی یه عضوی تو $W$ هست که برای یکی از $S$ هم تایید شده -- که دقیقاً همون چیزیه که JR میخواد.

EJR$\Rightarrow$PJR هم ساده‌ست: اگه EJR برقرار باشه یعنی یه $i^*\in S$ هست با $|A_{i^*}\cap W|\ge\ell$. چون $A_{i^*}\subseteq\bigcup_{i\in S}A_i$، پس $|(\bigcup_{i\in S}A_i)\cap W|\ge\ell$ هم برقراره -- که دقیقاً شرط PJR‌ه.

اون یکی، یعنی Core$\Rightarrow$EJR، یکم بیشتر کار میبره. فرض کن $W$ تو هسته باشه ولی EJR رو نقض کنه: یعنی یه $(S,\ell)$ هست با $|S|\ge\ell n/k$ و $|\bigcap_{i\in S}A_i|\ge\ell$، ولی همه‌ی $i\in S$ کمتر از $\ell$ نماینده تو $W$ دارن. حالا چون $|\bigcap_{i\in S}A_i|\ge\ell$، میشه یه زیرمجموعه $T^*$ از اون اشتراک با دقیقاً $\ell$ عضو انتخاب کرد. برای هر $i\in S$، چون $T^*\subseteq A_i$، داریم $|A_i\cap T^*|=\ell$، که از $\ell-1\ge|A_i\cap W|$ بیشتره. یعنی $|A_i\cap T^*|>|A_i\cap W|$. و از طرف شرط اندازه هم داریم $|T^*|=\ell\le\frac{|S|}{n}k$. پس $S$ با $T^*$ میتونه $W$ رو بلوکه کنه -- که با فرض هسته‌بودن $W$ در تضاده. پس هر کمیته‌ی هسته‌ای، EJR رو هم داره.

حالا چرا هرچی جلوتر میریم قوی‌تر میشه؟ به نظرم اینجوریه:
\begin{itemize}
\item JR فقط $\ell=1$ رو چک می‌کنه، و فقط میخواد \emph{یک نفر} تو گروه بزرگ، یه کاندیدای مشترک تو کمیته داشته باشه -- حتی اگه گروه خیلی بزرگ‌تر از حد لازم برای چند کرسی باشه.
\item PJR این رو به همه‌ی $\ell$ها تعمیم میده، ولی فقط اتحاد ترجیحات گروه رو چک می‌کنه. یعنی $\ell$ کرسی می‌تونه خیلی نامتوازن بین اعضا پخش بشه.
\item EJR این حفره رو می‌بنده: میخواد حداقل \emph{یک نفر مشخص} تو گروه، شخصاً $\ell$ کاندیدای موردعلاقه‌شو تو کمیته ببینه.
\item Core از این هم فراتر میره -- به جای اینکه فقط ساختار اشتراک $\ell$-تایی رو چک کنه، اجازه میده \emph{هر} گروه با \emph{هر} کمیته‌ی جایگزینِ متناسب اعتراض کنه، نه فقط گروه‌هایی که یه اشتراک بزرگ دارن.
\end{itemize}

برای اینکه مطمئن بشم این رابطه‌ها واقعاً یک‌طرفه‌ان (یعنی معکوسشون درست نیست)، با کد چند تا نمونه پیدا کردم:
\begin{itemize}
\item JR بدون PJR: $n=2, m=3, k=2$، هر دو نفر $A_i=\{a,c\}$، و $W=\{b,c\}$. JR برقراره ($c\in W$)، ولی تو $\ell=2$، $|\{a,c\}\cap W|=1<2$.
\item PJR بدون EJR: $n=2, m=4, k=2$، $A_1=\{a,b,c\}$، $A_2=\{b,c,d\}$، $W=\{a,d\}$. اتحادشون با $W$ اشتراک ۲تایی داره ولی هیچکدوم به تنهایی دو تا کاندیدا تو $W$ نداره.
\item EJR بدون Core: $n=3, m=3, k=2$، $A_1=\{a\}$، $A_2=A_3=\{a,b\}$، $W=\{b,c\}$. اینجا EJR کاملاً برقراره، ولی کل گروه با $T=\{a,b\}$ میتونه همه رو بهتر کنه -- چون $c$ اصلاً هیچکس تاییدش نکرده! این خودش نشون میده EJR حتی جلوی وجود یه کاندیدای کاملاً بی‌طرفدار تو کمیته رو هم نمی‌گیره.
\end{itemize}

\subsubsection*{پرسش ۲}

نه، $F_k^{AV}$ همیشه JR رو برآورده نمی‌کنه. این خودش همون چیزیه که تو مقدمه‌ی متن پروژه هم بهش اشاره شده: یه اکثریت بزرگ با ترجیح مشترک می‌تونه همه‌ی کرسی‌ها رو ببره.

نمونه: $n=10$، $m=6$، $k=3$. اکثریت (۶ نفر) دقیقاً $\{c_1,c_2,c_3\}$ رو تایید می‌کنن، اقلیت (۴ نفر) دقیقاً $\{c_4,c_5,c_6\}$ رو. امتیازها میشه $c_1,c_2,c_3$ هرکدوم ۶، و $c_4,c_5,c_6$ هرکدوم ۴. پس $F_3^{AV}$ دقیقاً $W=\{c_1,c_2,c_3\}$ رو انتخاب می‌کنه.

آستانه‌ی JR میشه $n/k=10/3\approx3.33$. اقلیت با ۴ نفر از این آستانه بالاتره و ترجیح مشترک هم داره، ولی هیچکدوم از اعضاش حتی یه نماینده هم تو کمیته نداره. یعنی JR واضحاً نقض شده.

\subsubsection*{پرسش ۳}

برای PAV: این یه نتیجه‌ی شناخته‌شده تو ادبیاته (Aziz و همکاران، ۲۰۱۷) که PAV همیشه EJR رو برآورده می‌کنه. من هم با جست‌وجوی تصادفی روی چند صد نمونه‌ی کوچیک هیچ نقضی پیدا نکردم که با این حرف جور درمیاد. فکر می‌کنم دلیلش وزن‌های هارمونیک $H_r$ باشه -- طوری طراحی شدن که اگه یه گروه بزرگ کافی با اشتراک کافی هیچ نماینده‌ی خوبی نداشته باشه، میشه با جایگزین کردن چندتا کاندیدای دیگه با کاندیداهای مشترک اون گروه، امتیاز کل PAV رو بالا برد -- شبیه همون استدلال «مبادله»ای که تو اثبات Core$\Rightarrow$EJR بالا زدم، فقط اینجا باید مراقب بود که امتیاز کل PAV هم بالا میره نه فقط مطلوبیت اون گروه. اثبات کاملش پیچیده‌تر از این حرفاست و تو منبعی که گفتم اومده، ولی ساختار کیفیش همینه.

برای CC اول با جست‌وجوی تصادفی هیچ نقض JR پیدا نکردم. ولی وقتی رفتم سراغ PJR، خیلی سریع یه نمونه پیدا شد: $n=2, m=3, k=2$، هر دو نفر دقیقاً $A_i=\{a,c\}$ رو تایید می‌کنن. چون $m-k=1$، هر کمیته‌ی ۲تایی از این ۳ کاندیدا حتماً یکی از $\{a,c\}$ رو داره -- یعنی هر سه کمیته‌ی ممکن امتیاز CC یکسان (۲) دارن. با شکست تساوی لغت‌نامه‌ای، $W=\{a,b\}$ برنده میشه. ولی تو $\ell=2$، اشتراکشون $\{a,c\}$ هست (اندازه ۲)، و $|S|=2\ge\ell n/k=2$، پس PJR باید $|\{a,c\}\cap W|\ge2$ رو تضمین کنه، ولی $\{a,c\}\cap\{a,b\}=\{a\}$، فقط ۱. نقض شده. و چون هیچکدوم به تنهایی هم به ۲ نمیرسه، EJR هم نقض میشه.

فکر می‌کنم فرقشون تو اینه که CC خیلی زود «سیر» میشه: مطلوبیت CC برای هر رأی‌دهنده صفر یا یکه -- به محض اینکه یه کاندیدای موردتاییدش تو کمیته باشه، برای CC دیگه فرقی نمی‌کنه چندتا کاندیدای دیگه‌ی همون آدم هم توش باشه. برای همین وقتی یه گروه هم‌رأی از قبل یه بار پوشش داده شده، CC هیچ انگیزه‌ای برای دادن کرسی دوم بهشون نداره، حتی اگه سهمشون ایجاب کنه. PAV چون وزناش هیچوقت کاملاً صفر نمیشن (فقط کمتر میشن) این مشکل رو نداره.

\subsubsection*{پرسش ۴}

مقایسه‌ی این سه قاعده یه تصویر جالب میده:

\begin{center}
\begin{tabular}{l|c|c}
قاعده & صداقت & بهترین تضمین نمایندگی \\
\hline
$F_k^{AV}$ & استراتژی برتر & هیچ‌کدوم (نه حتی JR) \\
CC & نه & JR \\
PAV & نه & EJR
\end{tabular}
\end{center}

الگویی که می‌بینم اینه: هرچی یه قاعده بیشتر تلاش می‌کنه نمایندگی رو بین گروه‌ها پخش کنه (با وابسته کردن امتیاز هر کاندیدا به کل ترکیب رأی هر نفر، نه فقط جمع ساده‌ی تک‌تک تاییدا)، همون چیزی که این نمایندگی رو ممکن می‌کنه (غیرجداپذیری) درِ سوءاستفاده رو هم باز می‌کنه. $F_k^{AV}$ چون امتیاز هر کاندیدا کاملاً مستقل حساب میشه صادقانه‌ست، ولی همین استقلال هیچ مکانیزمی برای حفظ سهم گروهی نداره.

از دید برگزارکننده، فکر می‌کنم این یه انتخاب ارزشیه: نبود تضمین نمایندگی تو $F_k^{AV}$ یه ریسک کوچیک و نادر نیست -- هر تقسیم‌بندی اکثریت/اقلیت با ترجیح جدا اونو فعال می‌کنه، همیشه اتفاق میفته. ولی نقطه‌ضعف PAV و CC که تو بخش قبل دیدیم به پیکربندی‌های نسبتاً خاص نیاز داشتن (برای پیدا کردن نمونه‌ی نقض $F_k^{AV}$ اولین تقسیم‌بندی ساده کافی بود، ولی برای PAV و CC هزاران نمونه‌ی تصادفی امتحان کردم تا یکی پیدا شد). پس اگه هدف اصلی جلوگیری از حذف کامل یه اقلیت بزرگه، به نظرم منطقی‌تره ریسک محدود PAV رو بپذیریم تا شکست تضمینیِ $F_k^{AV}$ رو.

\subsubsection*{پرسش ۵}

از همون نمونه‌ی پرسش ۲ استفاده می‌کنم: $n=10,k=3,m=6$، اکثریت $\{c_1,c_2,c_3\}$، اقلیت $\{c_4,c_5,c_6\}$، $W=\{c_1,c_2,c_3\}$.

اینکه پروفایل صادقانه تعادل نشه، مستقیم از پرسش ۱ بخش ۳ میاد (صداقت تحت $F_k^{AV}$ استراتژی برتره، پس برای این پروفایل هم برقراره). با کد هم چک کردم -- هیچ رأی‌دهنده‌ای با هیچ تغییری نمی‌تونه امتیازشو بالا ببره، چون فاصله‌ی امتیاز (۶ در برابر حداکثر ۵) بیشتر از چیزیه که یه رأی تنها بتونه جبران کنه.

ولی این کمیته تو هسته نیست. گروه بلوکه‌کننده رو کل اقلیت در نظر بگیر: $S=\{7,8,9,10\}$، $|S|=4$. سقف بودجه‌شون میشه $\frac{4}{10}\cdot3=1.2$، یعنی $T$ حداکثر یه کاندیدا می‌تونه داشته باشه. $T=\{c_4\}$ رو در نظر بگیر. برای هر $i\in S$:
\[
|A_i\cap T|=1 > 0=|A_i\cap W|.
\]
یعنی هر چهار نفر اقلیت با این پیشنهاد بهتر میشن، پس $S$ با $T$ کمیته رو بلوکه می‌کنه.

اینکه چرا این نشون میده انحراف فردی نبودن کافی نیست: تعادل نش فقط انحراف تک‌نفره رو چک می‌کنه -- آیا \emph{یه} رأی‌دهنده به تنهایی می‌تونه بهتر بشه؟ اینجا نه، چون رأی تک‌نفره در برابر فاصله ۶-به-۴ ناچیزه. ولی هسته یه سوال کاملاً متفاوت می‌پرسه: آیا یه \emph{گروه}، با جمع کردن قدرتشون، می‌تونن مشترکاً به نتیجه‌ی بهتری برسن؟ اینجا جوابش بله‌ست -- ۴۰٪ رأی‌دهنده‌ها که سهمشون از ۳ کرسی حدود ۱.۲ تاست، حتی یک کرسی هم نگرفتن. همین فرق بین «هیچکس تنهایی نمی‌تونه عوض کنه» و «یه گروه می‌تونست بهتر داشته باشه» دقیقاً همون چیزیه که هسته رو یه مفهوم جدا و قوی‌تر از تعادل نش می‌کنه.











% --- جایگزین کامل بخش پایانی پروژه ---
\section{پرسش نهایی: وجود هسته}

این بخش برای من از دو بخش قبل متفاوت بود. در سؤال‌های قبلی معمولاً می‌شد با یک مثال کوچک یا یک استدلال نسبتاً مستقیم به جواب رسید؛ اما اینجا خیلی زود به جایی رسیدم که معلوم نبود اصلاً باید دنبال «اثبات» بگردم یا «مثال نقض». بنابراین اول تلاش کردم مسئله را با محاسبه بررسی کنم و فقط وقتی دیدم در اندازه‌های کوچک هیچ اتفاقی نمی‌افتد، سراغ نتایج موجود در ادبیات رفتم.

\subsection{روند کار}

در شروع، هدفم پیدا کردن یک پروفایل بود که در آن هیچ کمیته‌ای در هسته نباشد. برای اندازه‌های کوچک، هم چند نمونه را دستی بررسی کردم و هم برنامه‌ای نوشتم که بتواند جست‌وجو را کامل انجام دهد. در این مرحله برخلاف بعضی از سؤال‌های قبلی، هیچ نمونه‌ی نقضی پیدا نشد. این نکته به‌تنهایی چیزی را ثابت نمی‌کرد، چون فضای جست‌وجو هنوز خیلی کوچک بود، اما حداقل نشان می‌داد که اگر قرار باشد هسته تهی شود، احتمالاً باید دنبال ساختار پیچیده‌تری بگردیم.

بعد از اینکه محاسبات را دوباره بررسی کردم و مطمئن شدم تعریف هسته و شرط بلوکه‌کردن را درست پیاده کرده‌ام، مسئله را در ادبیات دنبال کردم. اینجا نکته‌ی مهمی روشن شد: سؤال «آیا در هر انتخابات کمیته‌ای همیشه یک کمیته در هسته وجود دارد؟» هنوز به‌صورت کلی حل نشده است. این در واقع یک مسئله‌ی باز شناخته‌شده است که از کار Aziz و همکاران در سال ۲۰۱۷ باقی مانده و تا امروز پاسخ کامل برای آن پیدا نشده است. \cite{aziz2017core}

این نتیجه جهت جست‌وجوی من را عوض کرد. به‌جای اینکه صرفاً دنبال یک مثال نقض بگردم، بررسی کردم که چه حالت‌هایی تاکنون به‌طور قطعی حل شده‌اند و نتایج محاسبات خودم چقدر با آن‌ها سازگارند. همچنین برای بخش مربوط به PAV، بعد از پیدا کردن یک نمونه‌ی منتشرشده، آن را با محاسبه‌ی مستقل دوباره بررسی کردم تا مطمئن شوم نتیجه فقط حاصل نقل یک مثال از منبع نیست.

\subsection{جست‌وجوی محاسباتی خودم}

کدی که برای این بخش نوشتم، برای مقادیر کوچک، فضای پروفایل‌ها را کامل بررسی می‌کند. یعنی برای هر پروفایل ممکن، تمام کمیته‌های اندازه‌ی $k$ را در نظر می‌گیرم و بررسی می‌کنم که آیا حداقل یک کمیته شرط هسته را دارد یا نه. این جست‌وجو تصادفی نیست؛ بنابراین نبودن نمونه‌ی نقض در این محدوده واقعاً به این معنی است که تمام حالت‌های آن اندازه بررسی شده‌اند.

نتیجه‌ی محاسبات من این بود:

\begin{center}
\begin{tabular}{c|c|c|c}
$m$ & $n$ & $k$ & نتیجه \\
\hline
$2$--$4$ & $2$--$4$ & همه‌ی مقادیر معتبر & هسته همیشه غیرتهی بود \\
$5$ & $2$--$4$ & $1$--$4$ & باز هم هسته در همه‌ی پروفایل‌ها غیرتهی بود
\end{tabular}
\end{center}

این محاسبات البته جای یک اثبات کلی را نمی‌گیرند. چیزی که برای من مهم بود این بود که هیچ تناقضی بین نتیجه‌ی برنامه و نتایج نظری پیدا نکردم. در واقع، نتایج جدیدتر هم نشان می‌دهند که فضای حالت‌های «حل‌شده» بسیار بزرگ‌تر از چیزی است که با جست‌وجوی دستی می‌شود دید.

مهم‌ترین نتیجه در مقاله‌ی Dominik Peters در IJCAI 2025 به‌دست آمده است: هسته برای همه‌ی انتخابات با $k\le 8$ غیرتهی است، بدون هیچ محدودیتی روی $m$ و $n$. در همان کار، همچنین نشان داده می‌شود که وقتی $m\le15$ باشد، هسته برای هر $k\le m$ غیرتهی است. این نتایج با استفاده از برنامه‌ریزی خطی و جست‌وجوی محاسباتی به‌دست آمده‌اند. \cite{peters2025core}

نتیجه‌ی دیگری که بعد از آن منتشر شده، مسئله را از جهت تعداد رأی‌دهنده‌ها محدود می‌کند. Becker، Greger و Peters در سال ۲۰۲۶ نشان داده‌اند که اگر $n\le5$ باشد، هسته همیشه غیرتهی است. نتیجه حتی به حالتی با حداکثر پنج نوع رأی‌دهنده نیز تعمیم داده شده است؛ یعنی تعداد رأی‌دهنده‌ها می‌تواند بیشتر باشد، به شرطی که تعداد مجموعه‌های تأیید متمایز حداکثر پنج تا باشد. \cite{becker2026core}

در نتیجه، جست‌وجوی کوچک من بیشتر از اینکه یک نتیجه‌ی تازه باشد، برایم یک sanity check بود: در اندازه‌هایی که قابل بررسی کامل بودند، هیچ نمونه‌ی نقضی پیدا نشد و این با تصویری که از نتایج نظری موجود به دست می‌آید سازگار است.

\subsection{پرسش ۱ -- آیا همیشه هسته وجود دارد؟}

جواب کلی هنوز معلوم نیست. تا اینجا نه اثباتی وجود دارد که نشان دهد هر انتخابات approval-based یک کمیته‌ی هسته‌ای دارد و نه نمونه‌ای که نشان دهد ممکن است هسته تهی باشد. به همین دلیل، پاسخ دقیق به سؤال این نیست که «بله» یا «نه»؛ بلکه این است که مسئله در حالت کلی هنوز باز است.

با این حال، این به آن معنا نیست که اطلاعات کمی درباره‌ی آن داریم. برعکس، چند ناحیه‌ی مهم از مسئله کاملاً حل شده‌اند. از جمله:

\begin{itemize}
\item اگر $k\le8$ باشد، هسته همیشه غیرتهی است.
\item اگر $m\le15$ باشد، هسته همیشه غیرتهی است.
\item اگر $n\le5$ باشد، هسته همیشه غیرتهی است.
\item در حالت‌های خاصی مثل $k=m-1$ و $k=m$ نیز وجود هسته به‌طور کامل مشخص است.
\end{itemize}

این نتایج باعث می‌شوند که جست‌وجوی یک مثال نقض، اگر اصلاً چنین مثالی وجود داشته باشد، از محدوده‌ی ساده‌ی پارامترها بیرون برود. به‌خصوص نتیجه‌ی $k\le8$ نشان می‌دهد که برای پیدا کردن counterexample باید حداقل کمیته‌ای با اندازه‌ی ۹ را در نظر گرفت؛ نتیجه‌ی $m\le15$ هم همین محدودیت را از سمت تعداد کاندیداها ایجاد می‌کند. \cite{peters2025core}

نکته‌ای که در هنگام کار روی این سؤال برای من مهم شد، تفاوت بین «نبودن مثال نقض» و «وجود یک اثبات» بود. من با جست‌وجوی کامل برای اندازه‌های کوچک توانستم بگویم که در آن اندازه‌ها هسته همیشه وجود داشته است؛ اما از این واقعیت نمی‌شود نتیجه گرفت که این موضوع برای هر $n,m,k$ هم درست است. حتی نتایج خیلی قوی‌تر مقاله‌ها نیز هنوز این فاصله را به‌طور کامل پر نکرده‌اند. بنابراین در این قسمت بهتر است به‌جای حدس‌زدن، همان وضعیت فعلی مسئله را بپذیریم: \textbf{وجود همیشگی هسته در حالت کلی هنوز یک مسئله‌ی باز است.}

\subsection{پرسش ۲ -- آیا خروجی PAV همیشه در هسته است؟}

اینجا پاسخ روشن‌تر است: \textbf{خیر}. PAV با اینکه برای بسیاری از خواص نمایندگی قاعده‌ی بسیار قدرتمندی است، در حالت کلی نمی‌تواند هسته را تضمین کند. Peters و Skowron نیز نشان داده‌اند که هیچ قاعده‌ی welfarist نمی‌تواند هسته را به‌طور کلی تضمین کند و PAV در عوض تنها یک تقریب ضریب ۲ از هسته ارائه می‌دهد. \cite{peters2020welfarism}

برای دیدن این موضوع به شکل مشخص، از نمونه‌ی زیر استفاده می‌کنم. چهار رأی‌دهنده، ده کاندیدا و کمیته‌ای با اندازه‌ی هشت داریم:
\[
A_1=\{c_1,c_2,c_3\},\qquad
A_2=\{c_1,c_2,c_4\},\qquad
A_3=A_4=\{c_5,\dots,c_{10}\}.
\]

یکی از کمیته‌های برنده‌ی PAV در این مثال برابر است با
\[
W=\{c_1,c_2,c_5,c_6,c_7,c_8,c_9,c_{10}\}.
\]

امتیاز PAV این کمیته برابر با $79/10$ است و چند کمیته‌ی دیگر نیز به همین امتیاز می‌رسند. نکته‌ی اصلی برای سؤال ما این است که این کمیته در هسته نیست. گروه $S=\{1,2\}$ می‌تواند کمیته‌ی
\[
T=\{c_1,c_2,c_3,c_4\}
\]
را پیشنهاد کند. اندازه‌ی این کمیته دقیقاً در بودجه‌ی گروه قرار می‌گیرد، چون
\[
|T|=4\le\frac{2}{4}\cdot8=4.
\]
از طرف دیگر، برای هر دو رأی‌دهنده داریم
\[
|A_1\cap T|=3>2=|A_1\cap W|,
\qquad
|A_2\cap T|=3>2=|A_2\cap W|.
\]
پس هر دو نفر با این انحراف گروهی وضعیت بهتری پیدا می‌کنند و $S$، $W$ را بلوکه می‌کند.

این مثال یک نکته‌ی مهم دارد: نشان نمی‌دهد که \emph{هر} خروجی PAV در این پروفایل خارج از هسته است، چون بین چند کمیته تساوی وجود دارد و ممکن است یکی از آن‌ها در هسته باشد. اما برای رد کردن ادعای «PAV همیشه هسته را تضمین می‌کند» همین مثال کافی است: وجود یک خروجی PAV که در هسته نیست، برای نقض ادعا کافی است.

حتی از این هم قوی‌تر، محدودیت حاصل از welfarist بودن PAV یک مسئله‌ی ساختاری است و به یک مثال خاص وابسته نیست. بنابراین نمی‌توان با تغییر شیوه‌ی شکست تساوی، این مشکل را در حالت کلی برطرف کرد. \cite{peters2020welfarism}

در عین حال، نتایج Peters در سال ۲۰۲۵ نشان می‌دهند که PAV برای کمیته‌های کوچک رفتار بسیار بهتری دارد: برای $k\le7$ هر کمیته‌ای که PAV انتخاب کند هسته‌ای است، و برای $k=8$ دست‌کم یکی از کمیته‌های بهینه‌ی PAV در هسته قرار دارد. پس شکست عمومی PAV از همان اندازه‌های کوچک شروع نمی‌شود؛ مسئله زمانی جدی می‌شود که $k$ بزرگ‌تر شود. \cite{peters2025core}

برای من، این نتیجه کنار مسئله‌ی پرسش ۱ جالب است. از یک طرف هنوز نمی‌دانیم اصلاً در همه‌ی اندازه‌ها یک کمیته‌ی هسته‌ای وجود دارد یا نه. از طرف دیگر می‌دانیم یک قاعده‌ی بسیار طبیعی و مهم مثل PAV نمی‌تواند برای همه‌ی اندازه‌ها آن کمیته را تضمین کند.

\subsection{جمع‌بندی}

\begin{center}
\begin{tabular}{l|l}
پرسش & پاسخ \\
\hline
همیشه هسته وجود دارد؟ & در حالت کلی هنوز باز است \\
آیا برای بعضی حالت‌ها اثبات شده؟ & بله؛ از جمله $k\le8$، $m\le15$ و $n\le5$ \\
آیا PAV همیشه یک کمیته در هسته برمی‌گرداند؟ & نه، در حالت کلی چنین تضمینی ندارد
\end{tabular}
\end{center}

\subsection*{جمع‌بندی شخصی}

این بخش از پروژه در نهایت برخلاف چیزی که اول انتظار داشتم با یک قضیه‌ی تمیز و یک جواب قطعی تمام نشد. چیزی که پیدا کردم، بیشتر شبیه یک نقشه از چیزی بود که هنوز می‌دانیم و چیزی که هنوز نمی‌دانیم. از یک طرف، برای اندازه‌های مختلف نتایج قوی داریم؛ از طرف دیگر، دقیقاً در حالت کلی که سؤال اصلی مطرح می‌شود، هنوز یک علامت سؤال باقی مانده است.

به نظرم مهم‌ترین نکته‌ای که از این بخش یاد گرفتم، تفاوت بین سه چیز بود: \emph{وجود یک جواب پایدار}، \emph{پیدا کردن آن جواب با یک قاعده‌ی مشخص}، و \emph{قوی بودن آن قاعده نسبت به معیارهای دیگر}. این سه تا یکی نیستند. ممکن است هسته وجود داشته باشد ولی PAV آن را انتخاب نکند. ممکن است PAV در اندازه‌های کوچک هسته را تضمین کند ولی در اندازه‌های بزرگ نه. و حتی اگر در تعداد زیادی از نمونه‌های کوچک همه‌چیز خوب پیش برود، هنوز یک مثال نقض بزرگ‌تر می‌تواند کل ادعای کلی را خراب کند.

این دقیقاً همان چیزی است که در مقایسه‌ی PAV و $F_k^{AV}$ هم دیده می‌شد. $F_k^{AV}$ از نظر صداقت رفتاری بسیار تمیزی دارد، چون امتیاز هر کاندیدا جداگانه محاسبه می‌شود؛ اما همین استقلال باعث می‌شود هیچ تضمین جدی برای نمایندگی گروهی نداشته باشد. PAV تقریباً در جهت مخالف حرکت می‌کند: امتیاز رأی‌دهنده‌ها را طوری ترکیب می‌کند که گروه‌ها بتوانند نمایندگی متناسب‌تری بگیرند و در نتیجه EJR را تضمین می‌کند، اما همین وابستگی ساختاری باعث می‌شود هم انگیزه‌ی رأی‌دادن صادقانه ضعیف‌تر شود و هم تضمین هسته به‌طور کامل از دست برود. در ادبیات هم این تفاوت به‌صورت یک محدودیت کلی‌تر برای قواعد welfarist دیده می‌شود. \cite{peters2020welfarism}

در مورد خود سؤال وجود هسته، هنوز نمی‌توانم بگویم که پاسخ نهایی «بله» است. اما فکر می‌کنم بعد از نتایج ۲۰۲۵ و ۲۰۲۶، گفتن اینکه «هیچ چیز درباره‌ی آن نمی‌دانیم» هم دیگر درست نیست. الان چند محدوده‌ی بزرگ کاملاً حل شده‌اند و می‌دانیم اگر counterexampleای وجود داشته باشد، در یکی از محدوده‌های حل‌نشده قرار می‌گیرد. مهم‌تر از آن، کار جدید Becker، Greger و Peters نشان می‌دهد که حتی اضافه شدن محدودیت روی تعداد رأی‌دهنده‌ها می‌تواند وجود هسته را تضمین کند. \cite{becker2026core}

در نهایت، برای من ارزش این بخش دقیقاً در همین ناتمام‌بودن آن بود. در سؤال‌های قبلی می‌توانستم با یک محاسبه یا استدلال به «جواب» برسم؛ اینجا مجبور شدم بین چیزی که خودم با کامپیوتر بررسی کرده‌ام، چیزی که در یک مقاله ثابت شده، و چیزی که هنوز هیچ‌کس ثابت نکرده تفاوت بگذارم. به نظرم این برای یک پروژه‌ی نظریه‌ی بازی مهم‌تر از این بود که صرفاً یک جدول جواب داشته باشم. اگر قرار باشد یک نتیجه‌ی نهایی از کل پروژه بگیرم، آن نتیجه این است که انتخاب یک قاعده‌ی رأی‌گیری همیشه به معنی پیدا کردن یک راه‌حل بی‌نقص نیست؛ هر قاعده‌ای چیزی را بهتر می‌کند و در عوض در جای دیگری محدودیت ایجاد می‌کند. مسئله‌ی هسته فقط یکی از جاهایی است که این محدودیت خیلی واضح خودش را نشان می‌دهد.

\subsection*{منابع این بخش}

\begin{thebibliography}{9}

\bibitem{aziz2017core}
H. Aziz, S. Brill, V. Conitzer, E. Elkind, R. Freeman, and T. Walsh.
\newblock \emph{Justified Representation in Approval-Based Committee Voting}.
\newblock Social Choice and Welfare, 48, 461--485, 2017.

\bibitem{peters2020welfarism}
D. Peters and P. Skowron.
\newblock \emph{Proportionality and the Limits of Welfarism}.
\newblock Journal of Economic Theory, 2020.

\bibitem{peters2025core}
D. Peters.
\newblock \emph{The Core of Approval-Based Committee Elections with Few Seats}.
\newblock Proceedings of IJCAI 2025, 4014--4022, 2025.

\bibitem{becker2026core}
P. Becker, M. Greger, and D. Peters.
\newblock \emph{Core Existence in Approval-Based Committee Elections with up to Five Voter Types}.
\newblock arXiv:2605.06194, 2026.

\end{thebibliography}









\end{document}
